% 来源：OEBPS/Text/01_35.xhtml（OPF spine 顺序；仅第一册）
\cleardoublepage
\BookEntry[劳伦斯·A.克尼厄姆\q 《向格雷厄姆学思考，向巴菲特学投资》]{劳伦斯·A.克尼厄姆\\ 《向格雷厄姆学思考，向巴菲特学投资》}

\BookCover{images/书单20向格雷厄姆学思考，向巴菲特学投资.jpg}{[美] 劳伦斯·A.克尼厄姆著| 闾佳、侯君译| 机械工业出版社2011年版}

\begin{BookSummaryBox}
{\centering \textbf{《巴菲特致股东的信》作者对巴菲特理念的深入解读\\ 价值投资的经典著作}\par}
\end{BookSummaryBox}

\subsubsection*{这本书谈什么？}

这本书是《巴菲特致股东的信》作者对巴菲特理念的深入解读，也无疑是价值投资的经典著作。由格雷厄姆创立的价值投资理念，经过几十年的时间，由沃伦·巴菲特发扬光大，如今成为基本面分析的主流指导思想。
\par

本书作者劳伦斯·A.克尼厄姆是一位经验丰富的财会、金融领域的教授学者，对基本面的分析和巴菲特的实战操作都有深入研究。他在自己对价值投资与实践结合的基础上，撰写了这本在全世界投资人心中占有不朽地位的著作。本书能让你了解和真正认识价值投资理念，提升你的投资实战技巧，避免在未来的投资中走太多的弯路。
\par

\subsubsection*{芒格为什么谈起这本书？}

在“论基本的、普世的智慧”的演讲中，芒格谈到本杰明·格雷厄姆使用的投资方法：
\par

\begin{quotation}
格雷厄姆使用私人拥有价值的概念，也就是说，应该考虑整个企业出售给私人拥有者的话，能够卖多少钱。在很多情况下，是可以计算出来的。
\par

然后，你再把股价乘以股票的份数，如果你得到的结果是整个售价的三分之一或更少，他会说你买这样的股票是捡了大便宜。即使那是一家烂企业，管理者是个酗酒的老糊涂，每股的真实价值比你支付的价格高出那么多，这意味着你能得到各种各样的好处。你如果得到这么多额外的价值，用格雷厄姆的话来说，就拥有了巨大的安全边际。
\par
\end{quotation}

本杰明·格雷厄姆出生于伦敦，幼年时随全家迁居美国。他的父亲开过一家进口公司，但很快就关门大吉。虽然出身穷苦，格雷厄姆还是考上了哥伦比亚大学。从该校毕业后，他成为纽伯格·亨德森·劳伯公司的记录员。聪明才智让他很快脱颖而出，年仅25岁便成为该公司的合伙人。1929年的股市大崩盘差点让格雷厄姆破产，但他吸取了宝贵的投资教训。在20世纪30年代，格雷厄姆出版了一系列投资图书，后来都成了经典。这些书里面最著名的有《证券分析》和《聪明的投资者》。格雷厄姆提出了“内在价值”的概念，以及以该价值的折扣价购买股票的理念。
\par

\BookMinorHeading{《向格雷厄姆学思考，向巴菲特学投资》}

\subsubsection*{作者劳伦斯·A. 克尼厄姆如是说}

“投资的时候，保持简单的头脑，做显而易见符合常理的事情。巴菲特的建议是：难易程度并不决定你投资的成功与否，找一家能干得久的、有一定发展规律可循的、并且你知道他们是干什么的公司，按照你的常理投资。”
\par

“成大事的人需要有耐心。巴菲特说：你心里要想着十年，而不是十分钟。如果你没有做好把一只股票持上个十年的准备，你一开始就不应该买入这只股票。耐心加上远瞻的能力，是巴菲特成功的另一秘诀。”
\par

“畏首畏尾而‘不把鸡蛋放在一个篮子’里的人，确实很安全。但是如果你认定了一个公司有着很好的发展前景，为什么不变得有攻击性一些呢？巴菲特的建议：你的股票最好不要超过十只，避免资金过度分散；在投资的战场上，要有勇气。”
\par
